Binge drinking is bad, it leads to a lot of bad things, 
and some epi stats, which means we're important.

Binge drinking interventions often try to target...something.

But those somethings are often poorly operationalized or not at all, either through...idk...probably
self-report. Studies on self-report show
(what, i imagine they're pretty bad, but maybe they're actually surprisingly good). 
Wearable transdermal alcohol sensors have emerged as an alternative to self-report,
offering a passive continuous measure of alcohol consumption
\parencite{fairbairn2025,gunn2023}. It surpasses previous measures of intoxication (eg
breathalyzers) by reducing the burden traditionally involved in obtaining a measurement to (blakn)
 what were the findings of those feasability studies?  
But although we have all this data, we dont know what to do with it.
Raw transdermal alcohol concentration (TAC) signal
can be difficult to interpret and translate into meaningful drinking parameters
\parencite{gunn2023,kianersi2023}.

To bridge the gap between theory and methods---because this phrasing looks good for pamt---
, we propose a formal model for binge-drinking, which operationalizes
binged drinking by...---some number--- parameters characterizing...---the things we
characterize---. 

Drawing from evidence in...some binge-drinking applied literature...we conceptualize
binge-drinking as...the empirical phenomena that make up binge-drinking...
For example, TAC-derived features such as rise rate, peak concentration, and rise duration
have been shown to predict alcohol-induced blackouts in college drinkers
\parencite{richards2024}, and TAC signal processing models have been developed and validated
against self-reported and breath alcohol concentration data in similar populations
\parencite{kianersi2023}.
Our model uses ideas from...---insert literature, ie michaelis menten
or that one dynamical systems model of loneliness---
...to characterize---the big picture formulation ie between drinking time, within episode
duration, within episode drinking patterns---as a dynamical system. 
Compared to previous models, ours provide interpretable parameters with mechanistic meanings.
Furthermore, we are able to distinguish between the end of a drinking episode and 
a return to baseline TAC. 
We either
if its not estimatable, some estimation routine speed up or contribution 
if its actually fine, an in depth demonstration and evaluation of the model utility

simulation demonstrate the accuracy of our method/software.

We illustrate how these parameters could be used to characterize the effects of risk factors 
on binge-drinking in college students.

the remainder of this paper is organized as follows: ---big long run-on sentence---

\subsection{the model}

